主码扫描将相邻候选分配给相邻线程。原有 byte 布局已将同一子空间的编码放在连续地址；打包进一步让每个线程通过一次 word 读取取得四个子空间编码，减少读取操作数。

\subsection{地址映射与四字节打包}

在列表 $\ell$ 中，行 $r$ 的物理位置为
\[
 p=o_\ell+r,\qquad \operatorname{offset}(Z[m,p])=mN+p.
\]
固定 $m$ 时，$\Delta\operatorname{offset}=\Delta p$，所以相邻候选拥有相邻编码 byte；跨列表边界时可能失去这种相邻性。打包布局将同一候选的四个子空间编码组合起来：
\begin{equation}
 \mathsf{w}[j,p]=\sum_{t=0}^{3}2^{8t}Z[4j+t,p],\qquad
 \operatorname{offset}_{\rm bytes}(\mathsf{w}[j,p])=4(jN+p).
 \label{eq:5-packed-layout}
\end{equation}
移位和掩码可以恢复各个原始八位值：
\[
 (\mathsf w[j,p]\gg8t)\mathbin{\&}255=Z[4j+t,p].
\]
编码的高、低 nibble 随后索引 \cref{eq:5-lut-factorisation} 中的两张表。

\par\medskip\noindent\begin{minipage}{\linewidth}
\captionsetup{hypcap=false}
\centering
% One column is one independent lane; there is no cross-lane reduction.
\begin{tikzpicture}[font=\small\linespread{1}\selectfont,
 box/.style={draw=ThesisNavy!55,rounded corners=2pt,align=center,
 text width=3.1cm,minimum height=1.05cm,inner sep=4pt},
 arr/.style={-{Latex[length=2mm]},draw=ThesisNavy!80,thick}]
\node[font=\small\bfseries] at (5.7,1.05) {Fixed word group $j$; consecutive candidate positions};
\foreach \a in {0,1,2,3} {
 \node[box,fill=ThesisNavy!8] (w\a) at (3.8*\a,0) {lane $\a$: position $p+\a$\\$\mathsf w[j,p+\a]$};
 \node[box] (u\a) at (3.8*\a,-1.65) {Unpack own word\\4 code bytes};
 \node[box,fill=ThesisNavy!3] (l\a) at (3.8*\a,-3.30) {LUT lookups\\Accumulate score};
 \draw[arr] (w\a)--(u\a);\draw[arr] (u\a)--(l\a);}
\node[anchor=west,align=left] at (-1.55,-4.65)
 {Byte layout: 4 byte-load instructions per lane for 4 codes\\[3pt]
  Packed layout: 1 word-load instruction per lane for the same 4 codes};
\end{tikzpicture}

\captionof{figure}[候选并行的打包访问]{候选并行的打包访问.}
\label{fig:5-packed-access}
\end{minipage}\par\medskip
\cref{fig:5-packed-access} 中每个 lane 读取并解包自己的 word，各 lane 的候选分数彼此独立。构建将 $[M,N]$ bytes 转为 $[M/4,N]$ words，要求 $4\mid M$。直接把原数组转换成 word 类型会将不同候选的 bytes 混在一起，因此需要单独重打包。读写字节量及临时存储为
\[
 R=W=NM,\qquad S_{\rm old+new}=2NM.
\]
这次重打包不计入搜索时间。

\subsection{访存开销与评分一致性}

一次 word 读取替代每候选四次 byte 读取。对 32 个 lane，逻辑 word 窗口覆盖 $32\times4=128$ 个连续 bytes；实际事务数取决于对齐和活跃 lane。每候选载荷仍为 $M$ bytes，解包保留每个编码。两种被测试布局都按子空间组织候选，因而该对比没有单独分离转置或缓存效应。\Cref{sec:7-3} 测量其集成吞吐变化。
